% Банк упражнений — глава 4
% Глава определяется именем этого файла.

\begin{exo}[Интегрирование дифференциальной формы вдоль пути]{integrale-forme-differentielle}
\keywords{дифференциальная форма, точный дифференциал, критерий Шварца, криволинейный интеграл}

Рассмотрим две дифференциальные формы
\[
\omega_1 = y\,\mathrm{d}x + x\,\mathrm{d}y,
\qquad
\omega_2 = y\,\mathrm{d}x,
\]
и три пути, соединяющие начало координат $O(0,0)$ с точкой $M(1,1)$: путь
$\gamma_1$, состоящий из горизонтального отрезка $O\to(1,0)$, а затем
вертикального отрезка $(1,0)\to M$; путь $\gamma_2$, состоящий из
вертикального отрезка $O\to(0,1)$, а затем горизонтального отрезка
$(0,1)\to M$; и путь $\gamma_3$ — прямолинейный отрезок, непосредственно
соединяющий $O$ и $M$.

\begin{enumerate}
\item Покажите, что $\omega_1$ является точной формой, найдя функцию $f$, такую что $\omega_1 = \mathrm{d}f$.
\item Вычислите $\int_{\gamma_1}\omega_1$ и $\int_{\gamma_2}\omega_1$, параметризуя каждый отрезок. Получите результат повторно без вычислений.
\item Примените критерий Шварца к $\omega_2$. Какой вывод можно сделать?
\item Вычислите $\int_{\gamma_1}\omega_2$, $\int_{\gamma_2}\omega_2$ и $\int_{\gamma_3}\omega_2$. Прокомментируйте результат.
\end{enumerate}

\begin{indication}
На горизонтальном отрезке $\mathrm{d}y = 0$, а на вертикальном
$\mathrm{d}x = 0$: каждый интеграл сводится к обычному определённому интегралу.
\end{indication}

\begin{solution}
\textbf{Вопрос 1.} Подходит $f(x,y) = xy$, поскольку $\partial f/\partial x = y$ и $\partial f/\partial y = x$.

\textbf{Вопрос 2.} Для кривой, параметризованной как
$t\mapsto\bigl(x(t),y(t)\bigr)$, имеем
\[
\int_\gamma\omega_1
= \int \left[y(t)x'(t)+x(t)y'(t)\right],\mathrm{d}t.
\]

Путь $\gamma_1$ является объединением двух отрезков. На первом
$x(t)=t$ и $y(t)=0$, где $t\in[0,1]$; следовательно, $\mathrm{d}x=\mathrm{d}t$
и $\mathrm{d}y=0$. На втором $x(t)=1$ и $y(t)=t$, также при
$t\in[0,1]$; следовательно, $\mathrm{d}x=0$ и $\mathrm{d}y=\mathrm{d}t$. Поэтому
\[
\begin{aligned}
I_1
= \int_{\gamma_1}\omega_1
&= \int_0^1\left(0\times 1+t\times 0\right)\,\mathrm{d}t
 + \int_0^1\left(t\times 0+1\times 1\right)\,\mathrm{d}t \\
&= 0+\int_0^1\mathrm{d}t = 1.
\end{aligned}
\]

Для $\gamma_2$ первый отрезок параметризуется как $x(t)=0$, $y(t)=t$,
а второй — как $x(t)=t$, $y(t)=1$, причём в обоих случаях $t\in[0,1]$.
Отсюда
\[
\begin{aligned}
I_2
= \int_{\gamma_2}\omega_1
&= \int_0^1\left(t\times 0+0\times 1\right)\,\mathrm{d}t
 + \int_0^1\left(1\times 1+t\times 0\right)\,\mathrm{d}t \\
&= 0+\int_0^1\mathrm{d}t = 1.
\end{aligned}
\]
Оба значения совпадают. Это можно было предвидеть без вычислений: поскольку
$\omega_1=\mathrm{d}f$ при $f(x,y)=xy$, интеграл зависит только от концов пути,
\[
\int_\gamma\omega_1=f(M)-f(O)=f(1,1)-f(0,0)=1.
\]

\textbf{Вопрос 3.} Запишем $\omega_2=A(x,y)\,\mathrm{d}x+B(x,y)\,\mathrm{d}y$.
Здесь $A(x,y)=y$ и $B(x,y)=0$. Если бы форма была точной, критерий Шварца давал бы
\[
\frac{\partial A}{\partial y}=\frac{\partial B}{\partial x}.
\]
Но $\partial A/\partial y=1$, тогда как $\partial B/\partial x=0$.
Смешанные производные различны, поэтому $\omega_2$ не является точной формой.
Её интеграл может зависеть от пути между $O$ и $M$, что подтверждает следующий расчёт.

\textbf{Вопрос 4.} Поскольку $\omega_2=y\,\mathrm{d}x$, вклад в интеграл
может дать только изменение $x$ при ненулевом значении $y$.
Используя предыдущие параметризации, для $\gamma_1$ получаем
\[
\begin{aligned}
J_1
=\int_{\gamma_1}\omega_2
&=\int_0^1 0\times 1\,\mathrm{d}t
  +\int_0^1 t\times 0\,\mathrm{d}t \\
&=0.
\end{aligned}
\]
На $\gamma_2$ первый отрезок удовлетворяет $x(t)=0$, поэтому
$\mathrm{d}x=0$; на втором $x(t)=t$, $y(t)=1$ и
$\mathrm{d}x=\mathrm{d}t$. Следовательно,
\[
\begin{aligned}
J_2
=\int_{\gamma_2}\omega_2
&=\int_0^1 t\times 0\,\mathrm{d}t
  +\int_0^1 1\times 1\,\mathrm{d}t \\
&=1.
\end{aligned}
\]
Наконец, диагональный отрезок $\gamma_3$ параметризуется явно:
\[
x(t)=t,\qquad y(t)=t,\qquad t\in[0,1],
\]
откуда $\mathrm{d}x=\mathrm{d}t$ и $\mathrm{d}y=\mathrm{d}t$. Поэтому
\[
J_3
=\int_{\gamma_3}\omega_2
=\int_0^1 y(t)x'(t)\,\mathrm{d}t
=\int_0^1 t\,\mathrm{d}t
=\frac12.
\]
Три пути имеют одинаковые концы, но дают три различных значения:
$J_1=0$, $J_2=1$ и $J_3=1/2$. Именно такая зависимость от пути характерна
для неточной дифференциальной формы.
\end{solution}
\end{exo}

\begin{exo}[Одинаковое изменение внутренней энергии при трёх способах передачи]{meme-delta-u-trois-transferts}
\keywords{первый закон термодинамики, внутренняя энергия, теплота, механическая работа, электрическая работа, калориметрия, функция состояния}

Жидкость, её калориметр и небольшой погружённый резистор образуют замкнутую
систему, макроскопически находящуюся в покое. Сосуд жёсткий, а полная
теплоёмкость системы считается постоянной:
\[
C=2{,}00\ \mathrm{kJ\,K^{-1}}.
\]
Требуется перевести систему из равновесного состояния
$A$ при $T_A=20{,}0\,^\circ\mathrm{C}$ в равновесное состояние
$B$ при $T_B=25{,}0\,^\circ\mathrm{C}$. Примем, что
\[
U(B)-U(A)=C(T_B-T_A).
\]
Изменения макроскопических кинетической и потенциальной энергий равны нулю.
Всеми потерями энергии во внешнюю среду пренебрегаем.

Независимо выполняют три следующих протокола, каждый раз начиная с одного и того же состояния $A$:
\begin{itemize}
\item[Протокол 1] Систему приводят в тепловой контакт с горячим источником, не сообщая ей никакой работы.
\item[Протокол 2] Стенки теплоизолированы. Лопастная мешалка перемешивает жидкость благодаря падению груза массы $m$ с высоты $h=5{,}00$ м. В конце мешалка и жидкость покоятся, а вся механическая энергия, потерянная грузом, передаётся системе. Принять $g=9{,}81\ \mathrm{m\,s^{-2}}$.
\item[Протокол 3] Система получает теплоту $Q_3=4{,}00$ кДж, а недостающая энергия подводится к резистору электрическим способом. Получаемая электрическая мощность постоянна и равна $\mathcal P=300$ Вт.
\end{itemize}

\begin{enumerate}
\item Вычислите общее для трёх протоколов изменение внутренней энергии $\Delta U=U(B)-U(A)$.
\item Для каждого из первых двух протоколов определите теплоту $Q$ и работу $W$, полученные системой. Во втором протоколе вычислите необходимую массу $m$.
\item Для третьего протокола вычислите полученную электрическую работу и продолжительность питания резистора.
\item Все три процесса имеют одинаковые начальное и конечное состояния. Объясните, почему их значения $Q$ и $W$ всё же могут различаться. Содержит ли система «теплоту» или «работу» в конечном состоянии $B$?
\end{enumerate}

\begin{indication}
Примените $\Delta U=Q+W$ с банковским правилом знаков. Для мешалки
полученная работа равна уменьшению потенциальной энергии груза $mgh$.
Энергия, пересекающая границу по электрическим проводам, учитывается как электрическая работа.
\end{indication}

\begin{solution}
\textbf{Вопрос 1.} Разность температур равна $T_B-T_A=5{,}0$ К. Поэтому
\[
\Delta U=C(T_B-T_A)
=2{,}00\times5{,}0
=10{,}0\ \mathrm{kJ}.
\]
Это значение зависит только от двух равновесных состояний $A$ и $B$.

\textbf{Вопрос 2.} В первом протоколе система не получает работы:
$W_1=0$. Поэтому первый закон даёт
\[
Q_1=\Delta U=10{,}0\ \mathrm{kJ}.
\]
Теплота положительна, поскольку энергия поступает в систему.

Во втором протоколе стенки теплоизолированы, поэтому $Q_2=0$.
Всё изменение внутренней энергии обусловлено механической работой мешалки:
\[
W_2=\Delta U=10{,}0\ \mathrm{kJ}.
\]
Так как $W_2=mgh$, то
\[
m=\frac{W_2}{gh}
=\frac{10{,}0\times10^3}{9{,}81\times5{,}00}
\simeq 204\ \mathrm{kg}.
\]
Большой порядок величины напоминает, что даже небольшому повышению температуры
уже соответствует значительная механическая энергия.

\textbf{Вопрос 3.} Первый закон требует
\[
W_3=\Delta U-Q_3=10{,}0-4{,}00=6{,}00\ \mathrm{kJ}.
\]
Эта работа положительна: источник электрического питания передаёт энергию системе.
При $W_3=\mathcal P\,\Delta t$ получаем
\[
\Delta t=\frac{W_3}{\mathcal P}
=\frac{6{,}00\times10^3}{300}
=20{,}0\ \mathrm{s}.
\]

\textbf{Вопрос 4.} Три пути соответственно дают
\[
(Q_1,W_1)=(10{,}0,0)\ \mathrm{kJ},
\qquad
(Q_2,W_2)=(0,10{,}0)\ \mathrm{kJ},
\]
и
\[
(Q_3,W_3)=(4{,}00,6{,}00)\ \mathrm{kJ}.
\]
Во всех случаях сумма $Q+W$ равна $10{,}0$ кДж и приводит к одному и тому же
изменению $\Delta U$. Внутренняя энергия является функцией состояния, тогда как
теплота и работа описывают переносы энергии во время процесса. В конечном
состоянии $B$ система обладает внутренней энергией, но не содержит ни «теплоты», ни «работы».
\end{solution}
\end{exo}

\begin{exo}[Сжатие при постоянном внешнем давлении]{compression-pression-exterieure-constante}
\seoready{true}
\keywords{работа сил давления, постоянное внешнее давление, сжатие, расширение, расширение в вакуум}

Газ сжимают от объёма $V_A$ до объёма $V_B<V_A$ при постоянном внешнем давлении $P_0$.

\begin{enumerate}
\item Вычислите работу, полученную газом.
\item Газ возвращается от $V_B$ к $V_A$ против того же внешнего давления $P_0$. Вычислите полученную работу и сравните её с работой при сжатии.
\item Повторите рассмотрение расширения от $V_B$ к $V_A$, когда газ расширяется в вакуум. Объясните три полученных знака.
\end{enumerate}

\begin{indication}
Используйте $W=-\int P_{\mathrm{ext}}\,\mathrm{d}V$. Интегрировать следует именно
\emph{внешнее} давление, в том числе когда процесс не является квазистатическим.
\end{indication}

\begin{solution}
\textbf{Вопрос 1.} Внешнее давление постоянно, поэтому
\[
W_{A\to B}=-P_0(V_B-V_A)=P_0(V_A-V_B)>0.
\]
При сжатии газ получает работу.

\textbf{Вопрос 2.} Для обратного процесса
\[
W_{B\to A}=-P_0(V_A-V_B)<0.
\]
Эти две работы противоположны, поскольку процессы совершаются против одного и того же постоянного внешнего давления.

\textbf{Вопрос 3.} В вакууме $P_{\mathrm{ext}}=0$, поэтому
\[
W_{B\to A}^{\mathrm{vide}}=0.
\]
Следовательно, расширение не обязательно означает отрицательную работу: необходимо, чтобы газ действительно оттеснял внешнюю силу.
\end{solution}
\end{exo}
